\chapter{The Residue Refuses To Disappear}

\section{The antagonist and the treasure}

The first loop closed clean: a number sent around the machine's encoding came back itself. But that was the training
case, chosen because its answer was \emph{nothing left over}. This chapter introduces the character the whole book is
built around, the one that appears the moment a loop does not close clean --- the \emph{residue}. It is the antagonist,
because it is the thing the machine would most like to be rid of and cannot be. And it is the treasure, because
everything the instrument will ever read is read off it. The residue is the leftover that will not vanish, and its
refusal to vanish is the entire plot.

Begin with where a residue comes from, in the plainest terms this gauge has. Whenever the machine represents a value it
cannot hold exactly --- a real approached by a finite condition, a number too long for the code it is written in --- it
keeps a finite front and leaves a tail. That tail is the residue: the part of the value the representation could not
carry. In the training loop the tail was empty, because the value fit the code exactly. In every interesting case it is
not empty, and the machine faces a choice about the tail it cannot escape: carry it, or throw it away. The residue is
precisely the tail the machine must decide what to do with, and the chapter is about why there is only one honest thing
to do.

The residue is not noise, and this is the first thing to get right about it. Noise is what a channel adds; the residue
is what the machine's own finiteness leaves behind. It is structured, not random: it is exactly the information the
representation dropped, the difference between the value and the value-as-written. Information theory has a sharp name
for a value with no shorter description than itself --- it is \emph{incompressible}, and Chaitin~\cite{chaitin1974} and
Kolmogorov~\cite{kolmogorov1965} made the notion exact: most values cannot be squeezed below their own length, and the
part that resists squeezing is not a flaw in the compressor but a real property of the value. The residue is that
resistant part, surfaced. When the machine's code cannot shorten a value to fit, what will not fit is the incompressible
core, and that core does not disappear because the machine looked away from it.

In the two verbs the residue is what a funge cannot swallow. All chapter long the machine has funged --- pooled the
indistinguishable, gathered like into like --- and every funge that closes leaves nothing over. The residue is the
funge that does not close: the remainder the pooling cannot absorb, because it is precisely the part that refuses to be
made the same as anything already in the bin. The tange decided it apart, and no funge can un-decide it. So the residue
enters the story as the first thing the machine's two verbs cannot dispose of --- the leftover that is neither cleanly
told apart nor cleanly pooled away, and therefore the one thing the machine must learn to \emph{carry}.

\section{The lie of discarding}

The machine has a residue on its hands and two things it could do with it. It could carry the residue --- keep the tail,
account for it, let it ride along in the representation as an acknowledged remainder. Or it could discard it --- round
the value off, drop the tail, and report the clean front as if it were the whole. This section is about the second
option, and about why it is not a shortcut but a \emph{lie}. The honesty of the entire instrument turns on this one
refusal: the machine does not throw the residue away.

See exactly what discarding does. To drop the tail is to declare, silently, that the value and its truncation are the
same value --- to funge together two things the machine has already tanged apart. But they are not the same; the tail is
the difference between them, and the machine decided that difference into existence. To discard it is to assert a
sameness the machine's own procedure denied. That is the precise structure of the lie: not an error of arithmetic, but a
false report, a claim of \emph{equal} where the machine holds \emph{unequal}. A truncation reported as a value is a
distinction the machine made and then pretended it had not. The lossy step is dishonest not because it loses precision
--- every finite machine loses precision --- but because it loses precision and \emph{does not say so}.

There is a lawful version of losing information and an unlawful one, and the difference is the whole of this section's
ethics. Lossy compression, honestly done, announces its loss: it says \emph{I kept this much and dropped that much},
and the dropped part is named, bounded, accounted for. The data-processing inequality~\cite{cover2006} makes the price
formal --- information thrown away downstream cannot be recovered upstream, so a discard is irreversible and had better
be declared. The unlawful version is the silent one: the machine drops the tail and reports the front with no mark that
anything was dropped, so that a reader downstream takes the truncation for the value. The machine of this book is
forbidden the silent discard. When it cannot carry a tail exactly, it carries a \emph{bound} on the tail --- an honest
statement of how much it dropped --- and never a clean number pretending the tail was zero.

And here is the reason the discipline is not optional. The whole instrument is being built to report a number it cannot
compute exactly, and the only thing standing between an honest bracket and a crank's point is this refusal to discard.
A machine that drops its residues quietly can make any value look exact; it can round its way to a confident single
number and hide, in the dropped tails, all the imprecision that would have shown the number was never earned. The honest
machine cannot do this, by rule. It keeps every residue it cannot resolve, as a bound, so that when the final reading
comes it comes with its imprecision visible --- as the width of a bracket, not the false cleanliness of a point. The lie
of discarding is exactly the lie the ending refuses; the machine learns not to tell it here, on small residues, so that
it will not tell it at the end, on the one that matters.

In the two verbs the lie of discarding is a funge the machine has no right to perform. To drop the tail is to pool the
truncation with the value --- to declare them one bin --- when the tange has already ruled them two. The honest machine
refuses that funge: it keeps the residue tanged apart, held as its own thing, bounded and carried. Where the last
chapter's first loop closed a funge legitimately, because nothing was left over, discarding would close a funge
illegitimately, by throwing the leftover out the back. The machine's honesty is precisely its refusal to funge what it
has tanged apart and cannot honestly pool --- to carry the residue rather than lie it away.

\section{What the residue will become}

The machine has resolved to carry the residue rather than discard it. This closing section says why that resolution is
worth the trouble --- what the carried residue turns into, downstream, once the instrument is fully assembled. The
answer is the promise the rest of the book pays off: the residue, carried faithfully, is the thing every later reading
is a reading \emph{of}. It is not waste the machine tolerates. It is the raw material the whole instrument exists to
measure.

Hold the residue as this gauge sees it: an incompressible remainder, a bounded tail the representation could not carry,
kept honestly rather than dropped. That single object, the far chapters will show, is what the machine reads as a
\emph{cost} when it weighs its own work, as a \emph{code length} when it asks how much a value takes to describe, as a
\emph{bound} when it brackets a value it cannot pin. The residue is the common ancestor of every reading the
computational gauge takes. Nothing new will be introduced to be measured; the machine will only ever put this same
carried remainder under different apparatus and read what it says. The treasure was named an antagonist because the
machine would like to lose it; it is a treasure because it is, quite literally, all there is to measure.

The other gauges read the same carried residue in their own vocabularies, and it is worth naming that the readings
diverge, because the divergence is the point of there being four books. Where this gauge reads the residue as a cost or
a code length, a physical gauge reads it as mass, or as charge, or as a phase, or as the curvature of some surface; a
gauge that walks the code reads it as a particular quantity printed by a particular run. Those are other volumes'
names, and this book neither builds nor defends them --- there is no mass in this construction, no charge, no field,
only the carried remainder and the readings this gauge can take of it. The book names the physical readings once, to
mark that the same residue wears them, and then lets them go: this volume keeps the incompressible remainder and the
cost, the code length, and the bound it becomes under \emph{computational} apparatus.

That the one residue supports so many readings is not loose analogy; it is the reason the ending can be trusted, and
the book plants that reason here. If four gauges that share no vocabulary all turn out to be reading the same carried
remainder, then a bracket they all arrive at is not one gauge's artifact but a fact about the residue itself. The
machine carries the residue faithfully now, on small cases, precisely so that at the end four independent instruments
can each put it under their own apparatus and close on the same interval. The faithful carrying is what makes the
cross-check possible; discard the residue and there would be nothing for the four gauges to agree about.

In the two verbs the carried residue is the funge held open across the whole book. The machine tanged it apart --- a
remainder its representation could not absorb --- and rather than falsely funge it away, it keeps it, bagged as an
acknowledged leftover, to be read later. Every reading the instrument takes is a funge of this residue under some
apparatus: the cost is the residue pooled as a price, the bound is the residue pooled as an interval, the code length is
the residue pooled as a count of symbols. The residue refuses to disappear, and the machine's whole art is to keep it,
honestly, until it can be read --- which is the work the next part of the book, where the tower turns around, begins.
